% Første linje er dokumentklassen. Brug *altid* memoir!
\documentclass[a4paper,oneside,article]{memoir}

% Herunder indlæses forskellige pakker
\usepackage[utf8]{inputenc}  % Korrekt håndtering af æ, ø og å
\usepackage[T1]{fontenc}  % Korrekt håndtering af æ, ø og å
\usepackage{microtype}  % Typografisk magi! Giver bl.a. pænere orddeling
\usepackage{graphicx}  % Gør det muligt at indsætte billeder
\usepackage{amsmath}  % Giver adgang til uundværlige matematikting
\usepackage{mathtools}  % Retter ting i ams og giver ekstra muligheder
\usepackage[danish]{babel}  % Danske betegnelser og orddeling
\usepackage{lmodern}
\usepackage{alphabeta}
\usepackage{float} %Mere kontrol over billede-placering
\renewcommand{\danishhyphenmins}{22}  % Bedre dansk orddeling
\usepackage{siunitx} % Sientific notation for forskellige ting, med en regel om at den skal bruge cdot istedet for times. 
\sisetup{inter-unit-product = \cdot, 
        exponent-product=\cdot}
\usepackage{derivative} % Til at skrive differentialligninger 
\usepackage{xcolor} % For at få flere tekst farver 
\usepackage[colorlinks=true]{hyperref} % For at få refrences man kan trykke på, slet [] hvis man vil have kasser i stedet for farvet tekst.
\usepackage{pdfpages} % Til at indsætte andre pdf'er
\usepackage{subfig} % For at have to billeder ved siden af hinanden
\usepackage{unicode-math}

% Pænere toc indents 
\usepackage{tocloft}
\cftsetindents{section}{1em}{2em}
\cftsetindents{subsection}{4em}{0em}


\begin{document}
\author{Mikkel Moth Billing \\ 202307989}
\title{Statistical and Solidstate Physics}
\date{\today}
%\maketitle   % Indsætter overskriften
% Table of Contents indstillinger
\setcounter{tocdepth}{1}
\setcounter{secnumdepth}{1}
\hypersetup{linkcolor=black} % Laver tableofcontents sort.  
\tableofcontents %Indsætter ToC
\hypersetup{linkcolor=red} % og så alle andre hyperrefs røde. 

% Custom Commands
\newcommand{\ti}[1]{\cdot 10^{#1}} % Let: gange 10^x
\newcommand{\e}[1]{\emph{e}^{#1}} % Pænere e^x
\newcommand{\vm}[2]{\begin{pmatrix} #1 \\ #2 \end{pmatrix}}
\newcommand{\mat}[1]{\begin{pmatrix} #1 \end{pmatrix}}
\newcommand{\h}{\hbar}
\renewcommand{\v}[1]{\mathbfit{#1}}
\newcommand{\ep}{\epsilon_0}


% Lette paranteser 
\newcommand{\br}[2][1]{%
    \ifcase#1
    \or \left( #2 \right) %1 eller ingenting
    \or \left[ #2 \right] %2
    \or \left\langle #2 \right\rangle %3
    \or \left\{ #2 \right\} %4
    \or \left| #2 \right| %5
    \fi
}

% Til at sætter bileder ind i store dokumenter
\newcommand{\bil}[1]{
\begin{figure}[H]
    \centering
    \includegraphics[width=0.75\linewidth]{Billeder/#1.png}
    \label{fig:#1}
\end{figure}
}

% Til at Give equation numre 
\newcommand{\eq}[1]{\label{eq:#1}\tag{#1}}


% Lav ny pagestyle (kun odd fordi enkeltsidet)
\makepagestyle{Title}
\pagestyle{Title}
%\makeoddhead{Title}{\footnotesize\color{gray}{\theauthor}}{}{\textcolor{gray}{\thedate}}
\makeoddfoot{Title}{}{\footnotesize\color{gray}{\thepage{} / \thelastpage}}{}
% Også fod på siden med \maketitle
\makeoddfoot{plain}{}{\footnotesize\color{gray}{\thepage{} / \thelastpage}}{}
\newpage

% ----------Dokumentet Starter her! ------------
\chapter{Crystal Structures}

\section{General Description of Crystal Structures}
This section introduces the mathematical description of crystals.
\begin{itemize}
    \item \textbf{Lattice:} A lattice is a set of regularly spaced points defined by generating vectors. In three dimensions, a Bravais lattice is defined by vectors $\v{R}$ [p. 2]:
    \begin{align*}
        \v{R} = m\v{a}_1 + n\v{a}_2 + o\v{a}_3 \eq{1.2}
    \end{align*}
    Where $m, n, o$ are integers and $\v{a}_1, \v{a}_2, \v{a}_3$ are the primitive lattice vectors.
    \item \textbf{Basis:} The building block (atom or group of atoms) attached to every lattice point to form the crystal structure [p. 2].
    \begin{align*}
        \text{Lattice} + \text{Basis} = \text{Crystal Structure}
    \end{align*}
    \item \textbf{Unit Cell:} A volume that fills space without overlap when translated by lattice vectors.
    \begin{itemize}
        \item \textbf{Primitive Unit Cell:} Contains exactly one lattice point [p. 2].
        \item \textbf{Wigner-Seitz Cell:} A specific primitive unit cell defined as the region of space closer to one lattice point than any other [p. 2].
    \end{itemize}
\end{itemize}

\section{Some Important Crystal Structures}
Common structures found in metals and ionic solids.
\begin{itemize}
    \item \textbf{Cubic Structures:}
    \begin{itemize}
        \item \textbf{Simple Cubic (sc):} Rare due to low packing density (open structure) [p. 4].
        \item \textbf{Body-Centered Cubic (bcc):} One atom at corners and one in the center. Common for alkali metals [p. 4].
        \item \textbf{Face-Centered Cubic (fcc):} One atom at corners and one on each face. Highest possible packing density (74\%) [p. 4].
    \end{itemize}
    \item \textbf{Ionic Structures:}
    \begin{itemize}
        \item \textbf{CsCl Structure:} Two simple cubic lattices (Cs and Cl) shifted by half the body diagonal [p. 5].
        \item \textbf{NaCl Structure:} Two fcc lattices shifted by half the side length [p. 5].
    \end{itemize}
    \item \textbf{Close-Packed Structures:} Structures that maximize packing density.
    \begin{itemize}
        \item \textbf{Hexagonal Close-Packed (hcp):} Stacking sequence ABABAB... [p. 6].
        \item \textbf{Face-Centered Cubic (fcc):} Stacking sequence ABCABC... [p. 6].
    \end{itemize}
    Both have a coordination number (nearest neighbors) of 12 [p. 6].
\end{itemize}

\section{Crystal Structure Determination}
The primary method for determining structure is X-ray diffraction.

\subsection{X-Ray Diffraction (Bragg Theory)}
\begin{itemize}
    \item \textbf{Bragg Condition:} Constructive interference occurs when reflection from parallel lattice planes separated by distance $d$ satisfies [p. 8]:
    \begin{align*}
        n\lambda = 2d \sin \theta \eq{1.3}
    \end{align*}
    Where $\theta$ is the angle of incidence relative to the plane (not the normal).
    \item \textbf{Miller Indices $(hkl)$:} Integers describing lattice planes. They are derived from the reciprocal of the intercepts of the plane with the crystallographic axes [p. 8].
\end{itemize}

\subsection{General Diffraction Theory and the Reciprocal Lattice}
\begin{itemize}
    \item \textbf{Scattered Intensity:} In the kinematic approximation (single scattering), the intensity $I(\v{K})$ depends on the scattering vector $\v{K} = \v{k}' - \v{k}$ and the electron concentration $\rho(\v{r})$ [p. 11]:
    \begin{align*}
        I(\v{K}) \propto \left| \int_V \rho(\v{r}) \e{-i\v{K}\cdot\v{r}} dV \right|^2 \eq{1.11}
    \end{align*}
    \item \textbf{Reciprocal Lattice:} A set of vectors $\v{G}$ that satisfy $\e{i\v{G}\cdot\v{R}} = 1$ for all Bravais lattice vectors $\v{R}$ [p. 12]. The primitive reciprocal lattice vectors $\v{b}_1, \v{b}_2, \v{b}_3$ are defined such that [p. 12]:
    \begin{align*}
        \v{a}_i \cdot \v{b}_j = 2\pi\delta_{ij} \eq{1.17}
    \end{align*}
    \item \textbf{Laue Condition:} Constructive interference (diffraction) occurs only when the change in wave vector equals a reciprocal lattice vector [p. 14]:
    \begin{align*}
        \Delta \v{k} = \v{k}' - \v{k} = \v{G} \eq{1.25}
    \end{align*}
    This is equivalent to the Bragg condition.
    \item \textbf{Structure Factor:} The intensity of diffraction peaks depends on the basis of atoms within the unit cell. The integral over the unit cell splits into a sum over atoms (structure factor) and an atomic form factor [p. 15].
    \item \textbf{Ewald Construction:} A geometric visualization of the Laue condition. A sphere of radius $k = 2\pi/\lambda$ in reciprocal space determines which lattice points $\v{G}$ satisfy the diffraction condition [p. 16].
\end{itemize}

\subsection{Other Methods}
\begin{itemize}
    \item \textbf{Neutron Diffraction:} Neutrons interact with nuclei (stronger interaction with light atoms than X-rays) and have magnetic moments (useful for magnetic ordering) [p. 17].
    \item \textbf{Electron Diffraction:} Strong interaction with matter, limiting penetration depth. Suitable for surface studies (LEED) [p. 17].
\end{itemize}

\chapter{Bonding in Solids}

\section{Attractive and Repulsive Forces}
\begin{itemize}
    \item Bonding requires a balance between attractive forces (varying by bond type) and repulsive forces (Pauli exclusion principle at short range) [p. 23].
    \item \textbf{Interatomic Potential:} Often modeled as [p. 23]:
    \begin{align*}
        \phi(r) = \frac{A}{r^n} - \frac{B}{r^m} \eq{2.1}
    \end{align*}
    Where $n > m$ to ensure repulsion dominates at short distances.
\end{itemize}

\section{Ionic Bonding}
\begin{itemize}
    \item \textbf{Mechanism:} Electron transfer from electropositive to electronegative atoms, leading to electrostatic attraction between ions [p. 24].
    \item \textbf{Madelung Energy:} The total electrostatic energy of an ion in a crystal is the sum of interactions with all other ions [p. 24]:
    \begin{align*}
        E_{electrostatic} = -M_d \frac{e^2}{4\pi\ep a} \eq{2.2}
    \end{align*}
    Where $M_d$ is the \textbf{Madelung constant}, which depends only on the crystal structure, and $a$ is the nearest-neighbor distance.
    \item \textbf{Cohesive Energy:} The total energy required to separate the solid into isolated atoms. For ionic solids, it is the lattice energy (electrostatic + repulsive) minus the energy cost to create ions (ionization energy - electron affinity) [p. 25].
\end{itemize}

\section{Covalent Bonding}
\begin{itemize}
    \item \textbf{Mechanism:} Electrons are shared between neighboring atoms, forming highly directional bonds [p. 25].
    \item \textbf{Examples:}
    \begin{itemize}
        \item \textbf{Carbon:} Forms diamond ($sp^3$ hybridization, tetrahedral) and graphite/graphene ($sp^2$ hybridization, planar honeycomb) [p. 25].
        \item \textbf{Semiconductors:} Silicon (Si) and Germanium (Ge) also crystallize in the diamond structure [p. 7].
    \end{itemize}
    \item \textbf{Characteristics:} High stability and cohesive energies (several eV per atom). The directional nature determines the crystal structure more than packing density [p. 25].
\end{itemize}

\section{Metallic Bonding}
\begin{itemize}
    \item \textbf{Mechanism:} Valence electrons are delocalized (conduction electrons) and move freely between positive ion cores [p. 28].
    \item \textbf{Energy Gain:}
    \begin{itemize}
        \item \textbf{Kinetic Energy:} Delocalization reduces the curvature of the electron wave function ($\nabla^2\Psi$), lowering kinetic energy compared to localized atomic states [p. 28].
        \item \textbf{Potential Energy:} Electrons see an attractive potential from ions, but the interaction is complex due to electron-electron repulsion and the Pauli principle [p. 29].
    \end{itemize}
    \item \textbf{Structure:} Bonding is non-directional, favoring close-packed structures (fcc, hcp) to maximize overlap and delocalization [p. 29].
    \item \textbf{Transition Metals:} bonding is stronger due to a mix of delocalized $s/p$ electrons and more localized (covalent-like) $d$ electrons [p. 29].
\end{itemize}

\setcounter{section}{5}
\section{van der Waals Bonding}
\begin{itemize}
    \item \textbf{Mechanism:} Weak attraction arising from quantum mechanical fluctuations in the electron cloud. A transient dipole in one atom induces a dipole in a neighbor, leading to attraction [p. 29].
    \item \textbf{Characteristics:}
    \begin{itemize}
        \item Present in all solids but significant only when other bonding types are absent (e.g., noble gas crystals, interactions between graphite layers) [p. 30].
        \item Very weak binding energies ($\sim$ meV range) [p. 30].
        \item The interaction potential typically scales as $r^{-6}$ (Problem 2.4) [p. 31].
    \end{itemize}
\end{itemize}

\chapter{Mechanical Properties}

\section{Elastic Deformation}
This section covers the macroscopic and microscopic description of reversible deformation.

\subsection{Macroscopic Picture}
\begin{itemize}
    \item \textbf{Stress ($\sigma$):} Force per unit area, $\sigma = F/A$ (Pascal) [p. 33].
    \item \textbf{Strain ($\epsilon$):} Relative length change, $\epsilon = \Delta l / l$ (dimensionless) [p. 33].
    \item \textbf{Hooke's Law:} For small deformations, stress is linear with strain [p. 35]:
    \begin{align*}
        \sigma = Y\epsilon \eq{3.2}
    \end{align*}
    Where $Y$ is \textbf{Young's Modulus}.
    \item \textbf{Other Elastic Constants:}
    \begin{itemize}
        \item \textbf{Shear Modulus ($G$):} Relates shear stress $\tau$ to shear angle $\alpha$: $\tau = G\alpha$ [p. 35].
        \item \textbf{Bulk Modulus ($K$):} Relates hydrostatic pressure $p$ to relative volume change: $K = -V \frac{dp}{dV}$ [p. 36].
        \item \textbf{Poisson's Ratio ($\nu$):} Describes lateral contraction during tensile stress [p. 36]:
        \begin{align*}
            \frac{\Delta l_{transverse}}{l_{transverse}} = -\nu \frac{\Delta l_{longitudinal}}{l_{longitudinal}} \eq{3.6}
        \end{align*}
    \end{itemize}
\end{itemize}

\subsection{Microscopic Picture}
\begin{itemize}
    \item Elastic deformation corresponds to small displacements of atoms from their equilibrium positions in the interatomic potential $\phi(r)$ [p. 37].
    \item \textbf{Harmonic Approximation:} Expanding the potential around the equilibrium distance $a$ [p. 37]:
    \begin{align*}
        \phi(x) \approx \phi(a) + \frac{1}{2}\phi''(a)(x-a)^2 \eq{3.11}
    \end{align*}
    The quadratic term creates a linear restoring force (Hooke's law). The elastic constants are directly related to the curvature of the potential ($\phi''(a)$) [p. 37].
\end{itemize}

\chapter{Thermal Properties of the Lattice}

\section{Lattice Vibrations}
We treat atomic vibrations using the harmonic approximation (linear restoring forces).

\subsection{A Simple Harmonic Oscillator}
\begin{itemize}
    \item For a single atom of mass $M$ connected to a spring with force constant $\gamma$, the frequency is [p. 48]:
    \begin{align*}
        \omega = \sqrt{\frac{\gamma}{M}} \eq{4.2}
    \end{align*}
    \item \textbf{Amplitude:} Using the equipartition theorem ($\frac{1}{2}\gamma x_{max}^2 = k_B T$), the vibration amplitude at temperature $T$ is [p. 48]:
    \begin{align*}
        x_{max} = \sqrt{\frac{2k_B T}{\gamma}} \eq{4.5}
    \end{align*}
\end{itemize}

\subsection{An Infinite Chain of Atoms}
A 1D model for vibrations in a crystal.

\subsubsection{One Atom Per Unit Cell}
\begin{itemize}
    \item Consider a chain of atoms with mass $M$, spacing $a$, and spring constant $\gamma$. The equation of motion for the $n$'th atom is [p. 49]:
    \begin{align*}
        M\frac{d^2u_n}{dt^2} = -\gamma(2u_n - u_{n-1} - u_{n+1}) \eq{4.7}
    \end{align*}
    \item \textbf{Wave Solution:} We assume a solution of the form 
    \begin{align*}
        u_n = u \e{i(kan - \omega t)} \eq{4.8}.
    \end{align*}
    \item \textbf{Dispersion Relation:} Substituting the solution yields the relationship between frequency $\omega$ and wave vector $k$ [p. 49]:
    \begin{align*}
        \omega(k) = \sqrt{\frac{4\gamma}{M}} \left| \sin\left(\frac{ka}{2}\right) \right| \eq{4.10}
    \end{align*}
    \item \textbf{Long Wavelength Limit ($k \to 0$):} $\sin(x) \approx x$, so the dispersion becomes linear [p. 50]:
    \begin{align*}
        \omega(k) \approx \sqrt{\frac{\gamma}{M}} a k = v k \eq{4.11}
    \end{align*}
    Here $v$ is the \textbf{speed of sound}. Group velocity and phase velocity are equal.
\end{itemize}

\subsubsection{The First Brillouin Zone}
\begin{itemize}
    \item The dispersion relation $\omega(k)$ is periodic with period $2\pi/a$ (the reciprocal lattice vector).
    \item \textbf{First Brillouin Zone:} The range of unique $k$-values, defined as [p. 51]:
    \begin{align*}
        -\frac{\pi}{a} \le k < \frac{\pi}{a}
    \end{align*}
    \item Values of $k$ outside this zone describe the exact same physical motion of atoms as a corresponding $k$ inside the zone [p. 51].
    \item At the zone boundary ($k = \pm \pi/a$), the group velocity $d\omega/dk$ is zero, representing a standing wave where neighboring atoms move in opposite directions [p. 50].
\end{itemize}

\subsection{Two Atoms per Unit Cell}
\begin{itemize}
    \item Consider a 1D chain with two atoms per unit cell (masses $M_1$ and $M_2$, spacing $b$). The equations of motion are coupled, leading to a system of linear equations [p. 52].
    \item \textbf{Dispersion Relation:} Solving the determinant yields two branches of solutions [p. 52]:
    \begin{align*}
        \omega^2 = \gamma\left(\frac{1}{M_1} + \frac{1}{M_2}\right) \pm \gamma \sqrt{\left(\frac{1}{M_1} + \frac{1}{M_2}\right)^2 - \frac{4}{M_1 M_2} \sin^2\left(\frac{kb}{2}\right)} \eq{4.16}
    \end{align*}
    \item \textbf{Acoustic Branch ($-$):} $\omega \to 0$ as $k \to 0$. Atoms move in phase. Corresponds to sound waves [p. 53].
    \item \textbf{Optical Branch ($+$):} Finite $\omega$ at $k=0$. Atoms move out of phase. Can interact with electromagnetic waves (light) if the ions have opposite charges [p. 53].
\end{itemize}

\subsection{A Finite Chain of Atoms}
\begin{itemize}
    \item To model real solids, we use \textbf{Periodic Boundary Conditions} (Born-von Kármán): $u_{n+N} = u_n$ for a chain of $N$ atoms [p. 54].
    \item This discretizes the allowed wave vectors $k$:
    \begin{align*}
        k = \frac{2\pi}{Na} m \eq{4.20}
    \end{align*}
    Where $m$ is an integer. There are exactly $N$ allowed $k$-values in the first Brillouin zone, conserving the total number of degrees of freedom [p. 54].
\end{itemize}

\subsection{Quantized Vibrations, Phonons}
\begin{itemize}
    \item Treating lattice vibrations quantum mechanically leads to the concept of \textbf{phonons} (quanta of lattice vibration energy) [p. 56].
    \item The energy levels of a normal mode with frequency $\omega(k)$ are quantized as harmonic oscillators [p. 55]:
    \begin{align*}
        E_l(k) = \left(l + \frac{1}{2}\right)\h\omega(k) \eq{4.22}
    \end{align*}
    \item Phonons are bosons and follow Bose-Einstein statistics [p. 56].
\end{itemize}

\subsection{Three-Dimensional Solids}
\begin{itemize}
    \item Generalizing to 3D, the wave vector $\v{k}$ becomes three-dimensional.
    \item \textbf{Density of States $g(\omega)$:} The number of vibrational states per unit frequency interval. For a 3D linear dispersion $\omega = vk$, it is [p. 65]:
    \begin{align*}
        g(\omega) = \frac{V\omega^2}{2\pi^2 v^3} \eq{4.40}
    \end{align*}
\end{itemize}

\section{Heat Capacity of the Lattice}
This section deals with the lattice contribution to the specific heat $C_V = (\partial U / \partial T)_V$.

\subsection{Classical Theory (Dulong-Petit)}
\begin{itemize}
    \item According to the equipartition theorem, each atom has $3$ vibrational degrees of freedom with average energy $k_B T$ each (potential + kinetic) [p. 60].
    \item Total energy per mole is $3 N_A k_B T = 3RT$.
    \item \textbf{Heat Capacity:} $C = 3R \approx 24.9$ J/mol K. This matches experiment at high temperatures but fails at low temperatures where $C \to 0$ [p. 60].
\end{itemize}

\subsection{Einstein Model}
\begin{itemize}
    \item \textbf{Assumption:} All atoms vibrate independently with the same frequency $\omega_E$ [p. 62].
    \item Average energy is derived using Bose-Einstein statistics for the occupation number $\langle n \rangle$ [p. 62]:
    \begin{align*}
        \langle E \rangle = 3N_A \h\omega_E \left( \frac{1}{\e{\h\omega_E / k_B T} - 1} + \frac{1}{2} \right) \eq{4.31}
    \end{align*}
    \item \textbf{Result:} Explains the drop in $C$ at low $T$, but predicts an exponential decay ($\e{-\h\omega_E / k_B T}$), whereas experiments show a $T^3$ dependence [p. 62].
\end{itemize}

\subsection{Debye Model}
\begin{itemize}
    \item \textbf{Assumption:} Treat lattice vibrations as sound waves with linear dispersion $\omega = vk$ up to a cutoff frequency $\omega_D$ [p. 64].
    \item \textbf{Debye Frequency $\omega_D$:} Chosen such that the total number of modes equals $3N$ [p. 66]:
    \begin{align*}
        \omega_D^3 = 6\pi^2 \frac{N}{V} v^3 \eq{4.42}
    \end{align*}
    \item \textbf{Debye Temperature:} Defined as $\Theta_D = \h\omega_D / k_B$, for atoms it is related to their mass and bond force. [p. 66].
    \item \textbf{Low Temperature Limit ($T \ll \Theta_D$):} The heat capacity follows the \textbf{Debye $T^3$ law} [p. 66]:
    \begin{align*}
        C = \frac{12π^4}{5}Nk_B\br{\frac{T}{Θ_D}}^3 \eq{4.45}
    \end{align*}
    This agrees excellently with experimental data at low temperatures.
\end{itemize}

\section{Thermal Conductivity}
\begin{itemize}
    \item \textbf{Mechanism:} Heat is transported by phonons. We treat them as a gas of particles [p. 68].
    \item \textbf{Thermal Conductivity $\kappa$:} Given by the kinetic theory expression [p. 68]:
    \begin{align*}
        \kappa_p = \frac{1}{3} C v \lambda_p \eq{4.48}
    \end{align*}
    Where $C$ is the heat capacity per volume, $v$ is the phonon velocity (speed of sound), and $\lambda_p$ is the phonon mean free path.
    \item \textbf{Scattering Mechanisms:}
    \begin{itemize}
        \item \textbf{Low T:} Scattering by crystal defects/boundaries dominates. $\lambda_p$ is roughly constant (size of sample), so $\kappa \propto C \propto T^3$ [p. 68].
        \item \textbf{High T:} Phonon-phonon scattering dominates. The number of phonons increases with $T$, reducing $\lambda_p$. This causes $\kappa$ to decrease (typically as $1/T$) [p. 68].
    \end{itemize}
    \item \textbf{Umklapp Processes:} Anharmonic phonon-phonon scattering events where crystal momentum is not conserved (changed by a reciprocal lattice vector $\v{G}$). These are responsible for thermal resistance at high temperatures [p. 69].
\end{itemize}

\section{Thermal Expansion}
\begin{itemize}
    \item \textbf{Origin:} Thermal expansion arises from the \textbf{anharmonicity} of the interatomic potential (e.g., cubic terms in the Taylor expansion) [p. 70].
    \item In a purely harmonic potential, the mean separation between atoms would not change with temperature.
    \item As temperature increases, atoms vibrate with higher energy. due to the asymmetry of the potential (steeper repulsion, softer attraction), the average position shifts to larger separations [p. 71].
    \item \textbf{Coefficient of Thermal Expansion $\alpha$:} Defined as $\Delta l / l = \alpha \Delta T$. It typically vanishes as $T \to 0$ [p. 70].
\end{itemize}

\chapter{Electronic Properties of Metals: Classical Approach}

\section{Basic Assumptions of the Drude Model}
The Drude model (1900) applies kinetic gas theory to electrons in a metal.
\begin{itemize}
    \item \textbf{Independent Electron Approximation:} Electrons do not interact with each other [p. 77].
    \item \textbf{Free Electron Approximation:} Electrons do not interact with ions between collisions [p. 78].
    \item \textbf{Collisions:} Electrons collide with ion cores, which instantaneously thermalizes them [p. 78].
    \item \textbf{Relaxation Time ($\tau$):} The average time between collisions. The probability of a collision per unit time is $1/\tau$ [p. 78].
    \item \textbf{Electron Density ($n$):} Calculated assuming each atom contributes $Z_v$ valence electrons [p. 78].
\end{itemize}

\section{Results from the Drude Model}

\subsection{DC Electrical Conductivity}
\begin{itemize}
    \item \textbf{Drift Velocity:} In an electric field $\mathcal{E}$, electrons acquire an average drift velocity $\overline{v}$ [p. 79]:
    \begin{align*}
        \overline{v} = \frac{-e\mathcal{E}\tau}{m_e} \eq{5.4}
    \end{align*}
    \item \textbf{Ohm's Law:} The current density is $\v{j} = -ne\overline{v} = \sigma \mathcal{E}$. The conductivity $\sigma$ and resistivity $ρ$ is given by [p. 80]:
    \begin{align*}
        \sigma = \frac{ne^2\tau}{m_e} \eq{5.9} \\
        ρ = \frac{m_e}{ne^2τ} \eq{5.10}
    \end{align*}
    \item \textbf{Mobility ($\mu$):} Defined as the magnitude of drift velocity per unit field [p. 80]:
    \begin{align*}
        \mu = \frac{|\overline{v}|}{\mathcal{E}} = \frac{e\tau}{m_e} \eq{5.11}
    \end{align*}
\end{itemize}

\subsection{Hall Effect}
\begin{itemize}
    \item \textbf{Setup:} A magnetic field $B_z$ perpendicular to the current density $j_x$ exerts a Lorentz force, leading to a transverse electric field $\mathcal{E}_H$ (Hall field) [p. 81].
    \item \textbf{Hall Coefficient ($R_H$):} defined as $R_H = \mathcal{E}_H / (j_x B_z)$. The Drude model predicts [p. 82]:
    \begin{align*}
        R_H = \frac{-1}{ne} \eq{5.15}
    \end{align*}
    \item This provides a way to measure the carrier density $n$.
    \item \textbf{Failure:} For some metals (e.g., Al, Be), $R_H$ is positive, suggesting positive charge carriers, which contradicts the classical model [p. 82].
\end{itemize}

\subsection{Optical Reflectivity of Metals}
\begin{itemize}
    \item In an AC electric field $\mathcal{E}(t) \propto \e{-i\omega t}$, the electrons oscillate. The material's response is described by a frequency-dependent dielectric function $\epsilon(\omega)$ [p. 84]:
    \begin{align*}
        \epsilon(\omega) = 1 - \frac{\omega_p^2}{\omega^2} \eq{5.28}
    \end{align*}
    \item \textbf{Plasma Frequency ($\omega_p$):} The frequency where the metal transitions from reflecting to transparent [p. 84]:
    \begin{align*}
        \omega_p^2 = \frac{ne^2}{m_e \epsilon_0} \eq{5.29}
    \end{align*}
    \item For $\omega < \omega_p$, $\epsilon$ is negative and the metal reflects light. For $\omega > \omega_p$ (typically UV), it becomes transparent [p. 84].
\end{itemize}

\subsection{The Wiedemann-Franz Law}
\begin{itemize}
    \item Relates thermal conductivity $\kappa$ and electrical conductivity $\sigma$. The Drude model predicts [p. 85]:
    \begin{align*}
        \frac{\kappa}{\sigma} = L T = \frac{3}{2} \br{\frac{k_B}{e}}^2 T \eq{5.31}
    \end{align*}
    \item The constant $L$ is the \textbf{Lorenz number}. The prediction agrees well with experiment due to fortuitous cancellation of errors (velocity averaging) [p. 85].
\end{itemize}

\section{Shortcomings of the Drude Model}
\begin{itemize}
    \item \textbf{Electronic Heat Capacity:} Classical theory predicts $C_{el} = \frac{3}{2} n k_B$. Experiments show the electronic contribution is negligible at room temperature (solid obeys Dulong-Petit law) [p. 86].
    \item \textbf{Mean Free Path:} At low temperatures, the mean free path can be extremely long (mm range), which is inconsistent with electrons bouncing off fixed ions [p. 86].
    \item \textbf{ Alloys:} The strong resistivity increase in alloys is not well explained [p. 86].
\end{itemize}

\chapter{Electronic Properties of Solids: Quantum Mechanical Approach}

\section{The Idea of Energy Bands}
\begin{itemize}
    \item \textbf{Band Formation:} As atoms are brought together to form a solid, discrete atomic energy levels split and merge into continuous bands of allowed energies [p. 92].
    \item \textbf{Metals:} Characterized by a partially filled highest occupied band (e.g., Na, where the 3s band is half-full). Electrons can easily be excited into nearby empty states to conduct current [p. 93].
    \item \textbf{Insulators/Semiconductors:} Characterized by a completely filled valence band and a completely empty conduction band, separated by an energy gap ($E_g$). Electrons cannot easily gain energy [p. 94].
\end{itemize}

\section{Free Electron Model}
This model treats electrons as quantum particles in a box (the solid) with no interaction with the lattice ions ($U(\v{r})=0$).

\subsection{The Quantum Mechanical Eigenstates}
\begin{itemize}
    \item \textbf{Wave Functions:} Using periodic boundary conditions for a cube of volume $V$, the solutions are plane waves [p. 95]:
    \begin{align*}
        \psi_{\v{k}}(\v{r}) = \frac{1}{\sqrt{V}} \e{i\v{k}\cdot\v{r}} \eq{6.4}
    \end{align*}
    \item \textbf{Energies:}
    \begin{align*}
        E(\v{k}) = \frac{\h^2 k^2}{2m_e} \eq{6.6}
    \end{align*}
    \item \textbf{Fermi Energy ($E_F$):} At $T=0$, electrons fill states up to the highest energy $E_F$. In k-space, occupied states form a sphere with radius $k_F$ [p. 96]:
    \begin{align*}
        E_F = \frac{\h^2 k_F^2}{2m_e} = \frac{\h^2}{2m_e}(3\pi^2 n)^{2/3} \eq{6.11}
    \end{align*}
    Where $n$ is the electron density.
    \item \textbf{Density of States $g(E)$:} The number of states per energy interval increases with energy [p. 97]:
    \begin{align*}
        g(E) = \frac{V}{2\pi^2}\br{\frac{2m_e}{\h^2}}^{3/2} E^{1/2} \eq{6.13}
    \end{align*}
    \item \textbf{Finite Temperature:} Occupation is governed by the Fermi-Dirac distribution. Only electrons within $\sim k_B T$ of $E_F$ are affected ("soft zone") [p. 97].
\end{itemize}

\subsection{Electronic Heat Capacity}
\begin{itemize}
    \item \textbf{Pauli Principle:} Unlike classical gas particles, only a small fraction of electrons (those near $E_F$) can absorb thermal energy. Most are "frozen" in deep states [p. 99].
    \item \textbf{Result:} The electronic heat capacity is linear in temperature [p. 100]:
    \begin{align*}
        C_{el} = \frac{\pi^2}{3} k_B^2 g(E_F) T \eq{6.18}
    \end{align*}
    where $g(E_F)$ is the density of states at the Fermi Energy. 
    This explains why the electronic contribution is negligible at room temperature compared to the lattice ($3R$).
    This heat capacity is only for metals. 
\end{itemize}

\subsection{The Wiedemann-Franz Law}
\begin{itemize}
    \item The quantum model correctly predicts the ratio of thermal to electrical conductivity [p. 100]:
    \begin{align*}
        \frac{\kappa}{\sigma} = \frac{\pi^2}{3} \br{\frac{k_B}{e}}^2 T = L T \eq{6.19}
    \end{align*}
    The calculated Lorenz number $L$ agrees well with experiment (correcting Drude's fortuitous error).
\end{itemize}

\section{The General Form of the Electronic States}
\begin{itemize}
    \item \textbf{Bloch's Theorem:} For electrons in a periodic potential $U(\v{r})$, the eigenstates are the product of a plane wave and a lattice-periodic function $u_{\v{k}}(\v{r})$ [p. 104]:
    \begin{align*}
        \psi_{\v{k}}(\v{r}) = \e{i\v{k}\cdot\v{r}} u_{\v{k}}(\v{r}) \eq{6.26}
    \end{align*}
    \item \textbf{Properties:} 
    \begin{itemize}
        \item Electrons are delocalized over the entire crystal (mean free path can be very long) [p. 104].
        \item $\v{k}$ is the \textbf{crystal momentum}, a quantum number conserved modulo a reciprocal lattice vector $\v{G}$ [p. 109].
    \end{itemize}
\end{itemize}

\section{Nearly Free Electron Model}
\begin{itemize}
    \item \textbf{Assumption:} The lattice potential is a weak perturbation on free electrons.
    \item \textbf{Band Gaps:} Energy gaps open at the Brillouin zone boundaries ($k = n\pi/a$) due to Bragg reflection of the electron waves [p. 107].
    \item \textbf{Mechanism:} At the zone boundary, traveling waves form standing waves ($\psi(+) \sim \cos(\pi x/a)$ and $\psi(-) \sim \sin(\pi x/a)$). These have different probability densities relative to the ion cores, leading to different potential energies and thus a gap [p. 110].
\end{itemize}

\section{Tight-binding Model}
\begin{itemize}
    \item \textbf{Approach:} Starts with localized atomic orbitals and treats the solid as a giant molecule (Linear Combination of Atomic Orbitals) [p. 111].
    \item \textbf{Dispersion:} For a 1D chain of s-orbitals with nearest-neighbor overlap integral $\gamma$ [p. 114]:
    \begin{align*}
        E(k) = E_s - \beta - 2\gamma \cos(ka) \eq{6.62}
    \end{align*}
    \item \textbf{Bandwidth:} The width of the energy band is $4\gamma$. Stronger overlap (closer atoms) leads to wider bands [p. 114].
    \item \textbf{Wave Functions:} At $k=0$, orbitals add in phase (bonding-like). At $k=\pi/a$, they alternate sign (antibonding-like) [p. 115].
\end{itemize}

\section{Energy Bands in Real Solids}
\begin{itemize}
    \item \textbf{Metals:} Defined by having a Fermi energy that cuts through one or more bands. There is a finite density of states at $E_F$ (e.g., Al) [p. 117].
    \item \textbf{Semiconductors/Insulators:} The Fermi energy lies within an \textbf{absolute band gap}. The valence band is full, and the conduction band is empty (at $T=0$) [p. 118].
    \item \textbf{Direct:} A band gap is direct if the Valence Band Maximum (VBM) and Conducting Band Minimum (CBM) is at the same k-space (x-value). 
    \item \textbf{Filling Rule:} A single band can hold 2 electrons per unit cell. A solid with an odd number of valence electrons per unit cell is typically a metal (e.g., Na, Al) [p. 120].
    \item \textbf{Graphene:} A special semi-metal where the conduction and valence bands touch at the Dirac points (K points). It has a linear dispersion $E \propto k$ [p. 121].
\end{itemize}
\bil{band1}

\section{Transport Properties}
\begin{itemize}
    \item \textbf{Semiclassical Model:} Electrons are treated as wave packets moving with the \textbf{group velocity} [p. 122]:
    \begin{align*}
        \v{v}_g = \frac{1}{\h} \nabla_{\v{k}} E(\v{k}) \eq{6.47}
    \end{align*}
    \item \textbf{Dynamics:} An external electric field $\mathcal{E}$ changes the crystal momentum, not the velocity directly [p. 123]:
    \begin{align*}
        \h \frac{d\v{k}}{dt} = -e\mathcal{E} \eq{6.65}
    \end{align*}
    \item \textbf{Filled Bands:} A completely filled band carries no net current because for every electron with velocity $\v{v}$, there is one with $-\v{v}$ (sum over Brillouin zone is zero) [p. 124].
    \item \textbf{Effective Mass ($m^*$):} The response of an electron to a force depends on the band curvature [p. 125]:
    \begin{align*}
        m^* = \h^2 \br{\frac{d^2E}{dk^2}}^{-1} \eq{6.68}
    \end{align*}
    \begin{itemize}
        \item Near the bottom of a band, curvature is positive $\rightarrow$ electron-like ($m^* > 0$).
        \item Near the top of a band, curvature is negative $\rightarrow$ hole-like ($m^* < 0$).
    \end{itemize}
    \item \textbf{Holes:} Vacancies in a nearly filled band behave like positive charge carriers with positive mass [p. 124].
\end{itemize}

\chapter{Semiconductors}

\section{Intrinsic Semiconductors}
\begin{itemize}
    \item \textbf{Definition:} Materials with a band gap $E_g$ small enough (typically $< 3$ eV) to allow thermal excitation of electrons from the valence band (VB) to the conduction band (CB) [p. 131].
    \item \textbf{Chemical Potential ($\mu$):} At $T=0$, $\mu$ lies in the middle of the band gap. At finite $T$, charge neutrality ($n=p$) keeps it near the middle [p. 132].
    \begin{align*}
        μ = \frac{E_g}{2}+\frac{3}{4}k_bT\ln\br{\frac{m_h^*}{m_e^*}}. \eq{7.15}
    \end{align*}
    \item \textbf{Carrier Densities:}
    \begin{itemize}
        \item \textbf{Effective Mass Approximation:} The bands near the gap are modeled as parabolas with effective masses $m_e^*$ (electrons) and $m_h^*$ (holes) [p. 134].
        \item \textbf{Boltzmann Approximation:} If the chemical potential is far from the band edges ($|E-\mu| \gg k_B T$), the Fermi-Dirac distribution simplifies to a Boltzmann distribution [p. 136].
        \item \textbf{Electron Density ($n$):}
        \begin{align*}
            n = N_{C,eff} \e{-(E_g-\mu)/k_B T} \eq{7.11}
        \end{align*}
        Where $N_{C,eff}$ is the effective density of states in the CB.
        \item \textbf{Hole Density ($p$):}
        \begin{align*}
            p = N_{V,eff} \e{-\mu/k_B T} \eq{7.12}
        \end{align*}
        \item \textbf{Density of States for free electrons $g_C(E)$:}
        \begin{align*}
            g(E) = \frac{V}{2\pi^2}\br{\frac{2m_e^*}{\h^2}}^{3/2}\br{E-E_g}^{1/2} \eq{7.4}
        \end{align*}
        \item \textbf{Density of States for holes $g_V(E)$:}
        \begin{align*}
            g(E) = \frac{V}{2\pi^2}\br{\frac{2m_h^*}{\h^2}}^{3/2}\br{-E}^{1/2} \eq{7.6}
        \end{align*}
        For both $g$, $m^*$ is the effective mass, which is relative to the curve of the band. 
        
    \end{itemize}
    \item \textbf{Law of Mass Action:} The product $np$ is independent of the chemical potential (and thus doping) at a given temperature [p. 137]:
    \begin{align*}
        np = n_i^2 = 4\br{\frac{k_B T}{2\pi\h^2}}^3 (m_e^* m_h^*)^{3/2} \e{-E_g/k_B T} \eq{7.13}
    \end{align*}
\end{itemize}

\section{Doped Semiconductors}
Conductivity is controlled by adding impurities (dopants).
\begin{itemize}
    \item \textbf{n-type:} Donors (e.g., P in Si) have an extra valence electron. This electron is loosely bound (binding energy $\sim$ meV, large Bohr radius) and easily excited into the CB [p. 140].
    \item \textbf{p-type:} Acceptors (e.g., B in Si) capture an electron from the VB, creating a mobile hole [p. 141].
    \item \textbf{Temperature Regimes:}
    \begin{itemize}
        \item \textbf{Freeze-out (Low T):} Thermal energy is insufficient to ionize dopants [p. 141].
        \item \textbf{Extrinsic (Intermediate T):} All dopants are ionized. Carrier concentration is constant and determined by doping level ($n \approx N_d$) [p. 141].
        \item \textbf{Intrinsic (High T):} Thermal excitation across the gap dominates ($n_i \gg N_d$) [p. 141].
    \end{itemize}
\end{itemize}

\section{Conductivity of Semiconductors}
\begin{itemize}
    \item \textbf{Conductivity ($\sigma$):} Carried by both electrons and holes [p. 144]:
    \begin{align*}
        \sigma = e(n\mu_e + p\mu_h) \eq{7.19}
    \end{align*}
    \item \textbf{Temperature Dependence:} Unlike metals, $\sigma$ typically \textit{increases} with temperature because the exponential increase in carrier density ($n, p$) dominates the decrease in mobility ($\mu$) due to scattering [p. 144].
\end{itemize}

\section{Semiconductor Devices}

\subsection{The pn Junction}
\begin{itemize}
    \item \textbf{Formation:} Joining p-type and n-type materials creates a \textbf{depletion layer} where mobile carriers recombine, leaving fixed ionized dopants. This creates a built-in electric field and potential barrier [p. 145].
    \item \textbf{Currents:} Equilibrium is a balance between \textbf{diffusion current} (carriers moving due to concentration gradient) and \textbf{drift current} (minority carriers swept by the field) [p. 148].
    \item \textbf{Diode Behavior:}
    \begin{itemize}
        \item \textbf{Forward Bias:} External voltage opposes the built-in field, lowering the barrier. Diffusion current dominates $\rightarrow$ current flows exponentially [p. 149].
        \item \textbf{Reverse Bias:} External voltage adds to the built-in field, increasing the barrier. Only small leakage (drift) current flows [p. 149].
    \end{itemize}
\end{itemize}

\subsection{Other Devices}
\begin{itemize}
    \item \textbf{MOSFET:} A transistor where a gate voltage controls the channel conductivity. A sufficient gate voltage creates an \textbf{inversion layer}, allowing current to flow between source and drain (switching) [p. 150].
    \item \textbf{LED:} Radiative recombination of electrons and holes in a forward-biased pn-junction emits light with energy $h\nu \approx E_g$ [p. 151]. Requires a \textbf{direct band gap} (e.g., GaAs) for efficiency [p. 152].
    \item \textbf{Solar Cell:} Reverse process of LED. Light creates electron-hole pairs which are separated by the built-in field of the pn-junction, generating voltage [p. 152].
\end{itemize}

\chapter{Magnetism}

\section{Macroscopic Description}
\begin{itemize}
    \item \textbf{Magnetization ($\v{M}$):} Magnetic dipole moment per unit volume [p. 160].
    \item \textbf{Constitutive Relation:} In matter, the magnetic field $\v{B}$ is [p. 160]:
    \begin{align*}
        \v{B} = \mu_0 (\v{H} + \v{M}) \eq{8.3}
    \end{align*}
    \item \textbf{Susceptibility ($\chi_m$):} defined by $\mu_0 \v{M} = \chi_m \v{B}_0$ (linear regime).
    \begin{itemize}
        \item \textbf{Diamagnetic:} $\chi_m < 0$ (repelled by field) [p. 160].
        \item \textbf{Paramagnetic:} $\chi_m > 0$ (attracted by field) [p. 160].
    \end{itemize}
\end{itemize}

\section{Quantum Mechanical Description}
\begin{itemize}
    \item Classical physics (Bohr-van Leeuwen theorem) cannot explain magnetism; a quantum approach is required [p. 161].
    \item \textbf{Hamiltonian:} In a magnetic field $\v{B}$, the kinetic energy term changes ($\v{p} \to \v{p} - q\v{A}$). For a uniform field $\v{B} = B_0 \hat{z}$ [p. 162]:
    \begin{align*}
        H'_{kin} = \frac{\v{p}^2}{2m_e} + \frac{e}{2m_e}B_0(\v{r}\times\v{p})_z + \frac{e^2}{8m_e}B_0^2(x^2+y^2) \eq{8.12}
    \end{align*}
    \item \textbf{Paramagnetic Term:} The linear term ($\propto B_0$) involves angular momentum. It lowers energy when the magnetic moment aligns with the field [p. 163].
    \item \textbf{Diamagnetic Term:} The quadratic term ($\propto B_0^2$) is always positive, raising energy (Lenz's law effect on atomic scale) [p. 163].
    \item \textbf{Spin:} Electrons also have an intrinsic spin magnetic moment $\boldsymbol{\mu}_s = -g_e \mu_B \v{S}$ ($g_e \approx 2$) [p. 163].
\end{itemize}

\section{Paramagnetism and Diamagnetism in Atoms}
\begin{itemize}
    \item \textbf{Closed Shells:} Total angular momentum ($L$) and spin ($S$) are zero. Atoms with only closed shells are diamagnetic [p. 164].
    \item \textbf{Open Shells (Paramagnetism):} Electrons determine the total angular momentum $J$ using \textbf{Hund's Rules} to minimize energy (maximize $S$, then maximize $L$) [p. 165].
    \item \textbf{Magnetic Moment:} The total magnetic moment is related to $J$ by the Landé g-factor $g_J$ [p. 164]:
    \begin{align*}
        \mu_J = -g_J \mu_B m_J \eq{8.19}
    \end{align*}
    \item \textbf{Atomic Diamagnetism:} Arises from the quadratic term in the Hamiltonian. It is proportional to the mean square radius of the electron orbitals and is usually very weak [p. 166].
\end{itemize}

\section{Weak Magnetism in Solids}
Most solids exhibit weak magnetic responses categorized as diamagnetism or paramagnetism.

\subsection{Diamagnetic Contributions}
\begin{itemize}
    \item \textbf{Atomic Diamagnetism:} Induced currents in filled electron shells oppose the external field (Lenz's law). The susceptibility is negative and temperature-independent [p. 167]:
    \begin{align*}
        \chi_m = -\frac{\mu_0 Z N e^2}{6 V m_e} \langle r^2 \rangle \eq{8.24}
    \end{align*}
    where $Z$ is the number of electrons and $\langle r^2 \rangle$ is the mean square atomic radius.
    \item \textbf{Free Electron Diamagnetism (Landau):} Conduction electrons moving in helical paths also create a diamagnetic moment. It is small, typically $\chi_{Landau} = -\frac{1}{3} \chi_{Pauli}$ [p. 167].
\end{itemize}

\subsection{Paramagnetic Contributions}
\begin{itemize}
    \item \textbf{Curie Paramagnetism (Localized Moments):} Arises from atoms with unfilled shells (non-zero total angular momentum $J$).
    \begin{itemize}
        \item Magnetic moments align with the field against thermal disorder.
        \item \textbf{Curie's Law:} For $g_J \mu_B B_0 \ll k_B T$, susceptibility is inversely proportional to temperature [p. 169]:
        \begin{align*}
            \chi_m = \frac{C}{T} \quad \text{with} \quad C = \frac{\mu_0 N g_J^2 \mu_B^2 J(J+1)}{3 V k_B} \eq{8.28}
        \end{align*}
        \item \textbf{Hund's Rules} determine $L$, $S$, and $J$ for the ground state [p. 165].
        \item \textbf{Quenching:} In crystals (esp. 3d transition metals), the orbital moment $L$ is often "quenched" by the crystal field, so $J \approx S$ [p. 169].
    \end{itemize}
    \item \textbf{Pauli Paramagnetism (Free Electrons):}
    \begin{itemize}
        \item Conduction electrons with spin parallel to the field lower their energy, leading to a net imbalance in spin populations near the Fermi level [p. 170].
        \item Susceptibility is independent of temperature (determined by $g(E_F)$) [p. 171]:
        \begin{align*}
            \chi_m = \mu_0 \mu_B^2 g(E_F) \eq{8.32}
        \end{align*}
    \end{itemize}
\end{itemize}

\section{Magnetic Ordering}
Spontaneous ordering of magnetic moments occurs due to interactions between spins.

\subsection{Exchange Interaction}
\begin{itemize}
    \item Dipole-dipole interaction is too weak ($\sim 1$ K) to explain ordering at high $T$ [p. 172].
    \item \textbf{Exchange Interaction:} Originates from the interplay of the Pauli exclusion principle and Coulomb repulsion. For two electrons, the energy difference between symmetric (singlet) and antisymmetric (triplet) spatial states defines the exchange constant $X$ [p. 173].
\end{itemize}

\subsection{Heisenberg Model (Localized Spins)}
\begin{itemize}
    \item \textbf{Hamiltonian:} Describes interaction between nearest-neighbor spins [p. 175]:
    \begin{align*}
        H = -X \sum_{i, nn} \v{S}_i \cdot \v{S}_{nn} \eq{8.36}
    \end{align*}
    \item \textbf{Mean Field Approximation:} Spins feel an effective \textbf{Weiss Field} $\v{B}_W$ proportional to the magnetization $\v{M}$ [p. 175]:
    \begin{align*}
        \v{B}_W \propto X \v{M} \eq{8.40}
    \end{align*}
    \item \textbf{Ferromagnetism:} Spontaneous magnetization $M(T)$ exists below the \textbf{Curie Temperature} $\Theta_C$ [p. 176].
    \item \textbf{Curie-Weiss Law:} Above $\Theta_C$, the susceptibility follows [p. 177]:
    \begin{align*}
        \chi_m = \frac{C}{T - \Theta_C} \eq{8.44}
    \end{align*}
\end{itemize}

\subsection{Band Model (Itinerant Magnetism)}
\begin{itemize}
    \item Appropriate for 3d transition metals (Fe, Co, Ni) where magnetic electrons are delocalized.
    \item \textbf{Stoner Model:} Exchange interaction shifts the energy bands for spin-up and spin-down electrons relative to each other [p. 179].
    \item \textbf{Spontaneous Magnetization:} Occurs if the energy gained by shifting electrons to the majority spin band exceeds the kinetic energy cost. Favored by a high density of states $g(E_F)$ and strong exchange interaction (Stoner Criterion) [p. 179].
\end{itemize}

\subsection{Domains and Hysteresis}
\begin{itemize}
    \item \textbf{Domains:} To minimize external stray field energy, ferromagnets split into domains with different magnetization directions [p. 180].
    \item \textbf{Bloch Walls:} Transition regions between domains where spins rotate gradually [p. 180].
    \item \textbf{Hysteresis:} Magnetization under an external field involves irreversible domain wall motion (pinning by defects). Defines \textbf{Remanence} ($M_R$) and \textbf{Coercivity} ($B_C$) [p. 182].
\end{itemize}

\chapter{Dielectrics}

\section{Macroscopic Description}
\begin{itemize}
    \item \textbf{Polarization ($\v{P}$):} Dipole moment per unit volume. Linearly related to the electric field $\mathcal{E}$ via the susceptibility $\chi_e$ [p. 187]:
    \begin{align*}
        \v{P} = \chi_e \epsilon_0 \mathcal{E} \eq{9.1}
    \end{align*}
    \item \textbf{Dielectric Constant ($\epsilon$):} $\epsilon = 1 + \chi_e$. In a capacitor, the dielectric reduces the electric field to $\mathcal{E}_{ext}/\epsilon$ due to surface charges induced by polarization [p. 189].
\end{itemize}

\section{Microscopic Polarization}
Three main mechanisms contribute to $\v{P}$:
\begin{itemize}
    \item \textbf{Electronic Polarization:} Displacement of the electron cloud relative to the nucleus. Present in all atoms [p. 190].
    \item \textbf{Ionic Polarization:} Displacement of positive and negative ions in an ionic crystal lattice (e.g., NaCl) [p. 190].
    \item \textbf{Orientational Polarization:} Alignment of permanent molecular dipoles (e.g., in H$_2$O). Strong temperature dependence ($1/T$) [p. 190].
\end{itemize}

\section{The Local Field}
\begin{itemize}
    \item The field acting on a single atom, $\mathcal{E}_{loc}$, differs from the macroscopic average field $\mathcal{E}$.
    \item \textbf{Lorentz Local Field:} For a cubic lattice, the local field includes a correction from the polarization of the surroundings [p. 192]:
    \begin{align*}
        \mathcal{E}_{loc} = \mathcal{E} + \frac{\v{P}}{3\epsilon_0} = \frac{\epsilon + 2}{3} \mathcal{E} \eq{9.7}
    \end{align*}
    \item \textbf{Clausius-Mossotti Relation:} Connects the macroscopic $\epsilon$ to the microscopic atomic polarizability $\alpha$ [p. 192]:
    \begin{align*}
        \frac{\epsilon - 1}{\epsilon + 2} = \frac{N \alpha}{3 \epsilon_0 V} \eq{9.9}
    \end{align*}
\end{itemize}

\section{Frequency Dependence of the Dielectric Constant}
The dielectric function $\epsilon(\omega)$ describes the response to time-dependent fields (e.g., light).

\subsection{Excitation of Lattice Vibrations}
\begin{itemize}
    \item \textbf{Ionic Polarization:} At low frequencies, ions follow the field. At optical frequencies ($\sim 10^{14}$ Hz), ions are too heavy to respond, so only electronic polarization contributes [p. 193].
    \item \textbf{Oscillator Model:} Optical phonons ($\v{k} \approx 0$) are modeled as damped harmonic oscillators driven by the electric field [p. 193]:
    \begin{align*}
        \epsilon(\omega) = \epsilon_{opt} + \frac{Ne^2}{V\epsilon_0 M} \frac{1}{\omega_0^2 - \omega^2 - i\eta\omega} \eq{9.18}
    \end{align*}
    Where $\epsilon_{opt}$ is the high-frequency limit (electronic contribution), $\omega_0$ is the resonance frequency, and $\eta$ is damping.
    \item \textbf{Imaginary Part $\epsilon_i$:} Represents energy dissipation (absorption). It peaks at the resonance frequency $\omega_0$ [p. 195].
\end{itemize}

\subsection{Electronic Transitions}
\begin{itemize}
    \item For photon energies $h\nu > E_g$, electrons can be excited across the band gap. This leads to structure in $\epsilon(\omega)$ [p. 196].
    \item \textbf{Absorption:} Strong absorption occurs when the photon energy matches the energy difference between occupied (valence) and empty (conduction) states at the same $\v{k}$ (vertical transitions) [p. 197].
\end{itemize}

\chapter{Superconductivity}

\section{Basic Experimental Facts}
\begin{itemize}
    \item \textbf{Zero Resistivity:} Below a critical temperature $T_C$, resistivity drops to zero (upper limit $< 10^{-25}$ $\Omega$m). Current in a loop can persist for years [p. 204].
    \item \textbf{Critical Parameters:} Superconductivity is destroyed by high temperatures ($T > T_C$), high magnetic fields ($B > B_C$), or high current densities ($j > j_C$) [p. 206].
    \item \textbf{Meissner Effect:} A superconductor behaves as a perfect diamagnet ($\chi_m = -1$). It expels magnetic fields from its interior ($B=0$) [p. 207].
    \begin{itemize}
        \item This is distinct from a perfect conductor, which would trap flux rather than expel it [p. 208].
    \end{itemize}
    \item \textbf{Isotope Effect:} $T_C \propto M^{-1/2}$, where $M$ is the isotopic mass. This indicates that lattice vibrations (phonons) play a role in the mechanism [p. 209].
\end{itemize}

\section{Some Theoretical Aspects}

\subsection{Phenomenological Theory (London)}
\begin{itemize}
    \item \textbf{First London Equation:} Describes collisionless acceleration of superconducting electrons ($n_s$) [p. 210]:
    \begin{align*}
        \frac{\partial \v{j}}{\partial t} = \frac{n_s e^2}{m_s} \mathcal{E} \eq{10.3}
    \end{align*}
    \item \textbf{Second London Equation:} Describes the Meissner effect [p. 210]:
    \begin{align*}
        \nabla \times \v{j} = -\frac{n_s e^2}{m_s} \v{B} \eq{10.6}
    \end{align*}
    \item \textbf{Penetration Depth ($\lambda_L$):} Magnetic fields decay exponentially into the superconductor surface with a characteristic length $\lambda_L$ (typically $\sim 30-100$ nm) [p. 211]:
    \begin{align*}
        \lambda_L = \sqrt{\frac{m_s}{\mu_0 n_s e^2}} \eq{10.11}
    \end{align*}
\end{itemize}

\subsection{Microscopic BCS Theory}
Developed by Bardeen, Cooper, and Schrieffer (1957).
\begin{itemize}
    \item \textbf{Cooper Pairs:} A weak attractive interaction (mediated by phonons) causes electrons near the Fermi energy to form pairs with opposite spin and momentum ($\v{k}\uparrow, -\v{k}\downarrow$). These pairs behave like bosons [p. 214].
    \item \textbf{Mechanism:} An electron polarizes the lattice (attracts ions), creating a region of positive charge that attracts a second electron [p. 213].
    \item \textbf{BCS Ground State:} Cooper pairs condense into a coherent macroscopic quantum state.
    \item \textbf{Energy Gap ($2\Delta$):} An energy gap opens at the Fermi level in the single-particle excitation spectrum. Breaking a pair requires energy $2\Delta$ [p. 215].
    \item \textbf{Critical Temperature:} Predicted to be $k_B T_C \approx 1.13 \h\omega_D \e{-1/N(E_F)V}$, where $V$ is the interaction strength [p. 216].
\end{itemize}

\section{Experimental Detection of the Gap}
\begin{itemize}
    \item \textbf{Tunneling:} In a Superconductor-Insulator-Metal junction, current only flows when the voltage $eV$ exceeds the gap $\Delta$ [p. 218].
    \item \textbf{Optical Absorption:} Photons are absorbed only if $h\nu \ge 2\Delta$ (typically in the microwave/far-IR range) [p. 219].
    \item \textbf{Heat Capacity:} Below $T_C$, the electronic heat capacity drops exponentially ($C \propto \e{-\Delta/k_B T}$), confirming the existence of an energy gap [p. 220].
\end{itemize}

\setcounter{section}{4}
\section{Type I and Type II Superconductors}
\begin{itemize}
    \item \textbf{Type I:} Expels magnetic flux completely until $B_C$, then turns normal. Occurs when coherence length $\xi > \lambda_L$ (e.g., Pb, Sn) [p. 222].
    \item \textbf{Type II:} Has two critical fields, $B_{c1}$ and $B_{c2}$.
    \begin{itemize}
        \item For $B < B_{c1}$: Meissner state (complete expulsion).
        \item For $B_{c1} < B < B_{c2}$: \textbf{Vortex State}. Flux penetrates in quantized vortices (fluxoids) while the bulk remains superconducting.
        \item Occurs when $\xi < \lambda_L$ (e.g., Nb, alloys). High $B_{c2}$ makes them useful for magnets [p. 222].
    \end{itemize}
\end{itemize}

\section{High-Temperature Superconductivity}
\begin{itemize}
    \item Discovered in 1986 (Bednorz and Müller) in cuprate ceramics (e.g., YBCO).
    \item \textbf{Characteristics:} $T_C$ can exceed the boiling point of nitrogen (77 K) [p. 225].
    \item \textbf{Structure:} Layered perovskite structures with CuO$_2$ planes. Superconductivity is highly anisotropic (quasi-2D) [p. 226].
    \item \textbf{Mechanism:} Unconventional (non-BCS). Electron pairing still occurs, but likely not mediated by phonons (possibly magnetic excitations) [p. 226].
\end{itemize}

\end{document}